% Seminar 4 - Pagina 2. Fragment pentru
% % Seminar 4 - Pagina 2. Fragment pentru
% % Seminar 4 - Pagina 2. Fragment pentru
% % Seminar 4 - Pagina 2. Fragment pentru
% \input{seminars/pages/seminar4-pagina2-ro}; fără preambul sau pagină nouă.
% NOTE EDITORIALE:
% - Exercițiul 6 are două puncte. Demonstrația tratează separat
%   cazul unui monoid și cazul unui semigrup fără element neutru.
% - Sursa notează x*infinit = infinit*x = x, dar nu scrie separat
%   infinit*infinit = infinit; această egalitate este necesară pentru
%   ca operația să fie definită pe întregul M ca monoid.
% - Manuscrisul folosește generic Z_n. Egalitatea ulterioară
%   3-căciulă*2-căciulă = 1-căciulă*2-căciulă, cu 3 != 1,
%   este valabilă pentru n=4, nu pentru orice n. Păstrăm Z_n.
% - Exemplul final cu matrice este doar indicat în manuscris;
%   nu adăugăm o demonstrație care nu este scrisă acolo.
% - Prescurtările și textul foarte șters nu sunt reconstruite
%   dincolo de pasajele lizibile.

\subsection*{Seminar 4 - Pagina 2}

\noindent\underline{Exercițiul 6.}

\noindent\textbf{(a)} Să se arate că orice semigrup $S$ se poate
scufunda într-un monoid: există un monoid $M$ și un morfism injectiv
de semigrupuri
\[
  f:S\longrightarrow M,\qquad f(xy)=f(x)f(y).
\]

\noindent\textbf{(b)} Există și un monoid care nu se poate scufunda
într-un grup: există un monoid $M$ astfel încât nu există un grup $G$
și un morfism injectiv de monoizi $f:M\to G$.

\medskip
\noindent\underline{Demonstrație.}

\noindent\textbf{(a)} Fie $S$ un semigrup.

\smallskip
\noindent\textbf{I.} Dacă $S$ este monoid, luăm $M=S$ și
\[
  i:S\longrightarrow M=S,\qquad i(x)=x.
\]

\noindent\textbf{II.} Dacă $S$ este semigrup, dar nu este monoid,
fie $\infty\notin S$ și notăm
\[
  M\overset{\mathrm{not}}{=}S\cup\{\infty\}.
\]
Pentru $x,y\in M$, operația este definită în manuscris prin
\[
\begin{gathered}
  x\circ y\ \overset{\mathrm{def}}{=}\ xy
     \quad(x,y\in S),\\
  \infty\circ x\ \overset{\mathrm{def}}{=}\
     x\circ\infty\ \overset{\mathrm{def}}{=}\ x
     \quad(x\in S).
\end{gathered}
\]
Atunci $(M,\circ,\infty)$ este monoid. Fie
\[
  f:S\longrightarrow S\cup\{\infty\}=M,\qquad f(x)=x;
\]
$f$ este morfism injectiv de semigrupuri.

\medskip
\noindent\textbf{(b)} $(\mathbb Z_n,\cdot)$ este monoid, unde
\[
  \mathbb Z_n=\{\widehat{0},\widehat{1},\ldots,
       \widehat{n-1}\}.
\]
Presupunem că există un grup $G$ și un morfism injectiv de monoizi
$f:\mathbb Z_n\to G$, astfel încât
\[
  f(xy)=f(x)f(y),\qquad f(\widehat{1})=e.
\]
Argumentul prin simplificare din manuscris este
\[
\begin{aligned}
  f(\widehat{3})f(\widehat{2})
    &=f(\widehat{1})f(\widehat{2})\\
  &\Longrightarrow f(\widehat{3})=f(\widehat{1})\\
  &\overset{f\,\mathrm{inj.}}{\Longrightarrow}
      \widehat{3}=\widehat{1}\qquad\text{(fals).}
\end{aligned}
\]

\noindent La sfârșitul paginii mai este indicat monoidul de matrice
$(M_n(K),\cdot)$, cu $A,B,C\in M_n(K)$ astfel încât
\[
  AB=CB,\qquad A\neq C.
\]
; fără preambul sau pagină nouă.
% NOTE EDITORIALE:
% - Exercițiul 6 are două puncte. Demonstrația tratează separat
%   cazul unui monoid și cazul unui semigrup fără element neutru.
% - Sursa notează x*infinit = infinit*x = x, dar nu scrie separat
%   infinit*infinit = infinit; această egalitate este necesară pentru
%   ca operația să fie definită pe întregul M ca monoid.
% - Manuscrisul folosește generic Z_n. Egalitatea ulterioară
%   3-căciulă*2-căciulă = 1-căciulă*2-căciulă, cu 3 != 1,
%   este valabilă pentru n=4, nu pentru orice n. Păstrăm Z_n.
% - Exemplul final cu matrice este doar indicat în manuscris;
%   nu adăugăm o demonstrație care nu este scrisă acolo.
% - Prescurtările și textul foarte șters nu sunt reconstruite
%   dincolo de pasajele lizibile.

\subsection*{Seminar 4 - Pagina 2}

\noindent\underline{Exercițiul 6.}

\noindent\textbf{(a)} Să se arate că orice semigrup $S$ se poate
scufunda într-un monoid: există un monoid $M$ și un morfism injectiv
de semigrupuri
\[
  f:S\longrightarrow M,\qquad f(xy)=f(x)f(y).
\]

\noindent\textbf{(b)} Există și un monoid care nu se poate scufunda
într-un grup: există un monoid $M$ astfel încât nu există un grup $G$
și un morfism injectiv de monoizi $f:M\to G$.

\medskip
\noindent\underline{Demonstrație.}

\noindent\textbf{(a)} Fie $S$ un semigrup.

\smallskip
\noindent\textbf{I.} Dacă $S$ este monoid, luăm $M=S$ și
\[
  i:S\longrightarrow M=S,\qquad i(x)=x.
\]

\noindent\textbf{II.} Dacă $S$ este semigrup, dar nu este monoid,
fie $\infty\notin S$ și notăm
\[
  M\overset{\mathrm{not}}{=}S\cup\{\infty\}.
\]
Pentru $x,y\in M$, operația este definită în manuscris prin
\[
\begin{gathered}
  x\circ y\ \overset{\mathrm{def}}{=}\ xy
     \quad(x,y\in S),\\
  \infty\circ x\ \overset{\mathrm{def}}{=}\
     x\circ\infty\ \overset{\mathrm{def}}{=}\ x
     \quad(x\in S).
\end{gathered}
\]
Atunci $(M,\circ,\infty)$ este monoid. Fie
\[
  f:S\longrightarrow S\cup\{\infty\}=M,\qquad f(x)=x;
\]
$f$ este morfism injectiv de semigrupuri.

\medskip
\noindent\textbf{(b)} $(\mathbb Z_n,\cdot)$ este monoid, unde
\[
  \mathbb Z_n=\{\widehat{0},\widehat{1},\ldots,
       \widehat{n-1}\}.
\]
Presupunem că există un grup $G$ și un morfism injectiv de monoizi
$f:\mathbb Z_n\to G$, astfel încât
\[
  f(xy)=f(x)f(y),\qquad f(\widehat{1})=e.
\]
Argumentul prin simplificare din manuscris este
\[
\begin{aligned}
  f(\widehat{3})f(\widehat{2})
    &=f(\widehat{1})f(\widehat{2})\\
  &\Longrightarrow f(\widehat{3})=f(\widehat{1})\\
  &\overset{f\,\mathrm{inj.}}{\Longrightarrow}
      \widehat{3}=\widehat{1}\qquad\text{(fals).}
\end{aligned}
\]

\noindent La sfârșitul paginii mai este indicat monoidul de matrice
$(M_n(K),\cdot)$, cu $A,B,C\in M_n(K)$ astfel încât
\[
  AB=CB,\qquad A\neq C.
\]
; fără preambul sau pagină nouă.
% NOTE EDITORIALE:
% - Exercițiul 6 are două puncte. Demonstrația tratează separat
%   cazul unui monoid și cazul unui semigrup fără element neutru.
% - Sursa notează x*infinit = infinit*x = x, dar nu scrie separat
%   infinit*infinit = infinit; această egalitate este necesară pentru
%   ca operația să fie definită pe întregul M ca monoid.
% - Manuscrisul folosește generic Z_n. Egalitatea ulterioară
%   3-căciulă*2-căciulă = 1-căciulă*2-căciulă, cu 3 != 1,
%   este valabilă pentru n=4, nu pentru orice n. Păstrăm Z_n.
% - Exemplul final cu matrice este doar indicat în manuscris;
%   nu adăugăm o demonstrație care nu este scrisă acolo.
% - Prescurtările și textul foarte șters nu sunt reconstruite
%   dincolo de pasajele lizibile.

\subsection*{Seminar 4 - Pagina 2}

\noindent\underline{Exercițiul 6.}

\noindent\textbf{(a)} Să se arate că orice semigrup $S$ se poate
scufunda într-un monoid: există un monoid $M$ și un morfism injectiv
de semigrupuri
\[
  f:S\longrightarrow M,\qquad f(xy)=f(x)f(y).
\]

\noindent\textbf{(b)} Există și un monoid care nu se poate scufunda
într-un grup: există un monoid $M$ astfel încât nu există un grup $G$
și un morfism injectiv de monoizi $f:M\to G$.

\medskip
\noindent\underline{Demonstrație.}

\noindent\textbf{(a)} Fie $S$ un semigrup.

\smallskip
\noindent\textbf{I.} Dacă $S$ este monoid, luăm $M=S$ și
\[
  i:S\longrightarrow M=S,\qquad i(x)=x.
\]

\noindent\textbf{II.} Dacă $S$ este semigrup, dar nu este monoid,
fie $\infty\notin S$ și notăm
\[
  M\overset{\mathrm{not}}{=}S\cup\{\infty\}.
\]
Pentru $x,y\in M$, operația este definită în manuscris prin
\[
\begin{gathered}
  x\circ y\ \overset{\mathrm{def}}{=}\ xy
     \quad(x,y\in S),\\
  \infty\circ x\ \overset{\mathrm{def}}{=}\
     x\circ\infty\ \overset{\mathrm{def}}{=}\ x
     \quad(x\in S).
\end{gathered}
\]
Atunci $(M,\circ,\infty)$ este monoid. Fie
\[
  f:S\longrightarrow S\cup\{\infty\}=M,\qquad f(x)=x;
\]
$f$ este morfism injectiv de semigrupuri.

\medskip
\noindent\textbf{(b)} $(\mathbb Z_n,\cdot)$ este monoid, unde
\[
  \mathbb Z_n=\{\widehat{0},\widehat{1},\ldots,
       \widehat{n-1}\}.
\]
Presupunem că există un grup $G$ și un morfism injectiv de monoizi
$f:\mathbb Z_n\to G$, astfel încât
\[
  f(xy)=f(x)f(y),\qquad f(\widehat{1})=e.
\]
Argumentul prin simplificare din manuscris este
\[
\begin{aligned}
  f(\widehat{3})f(\widehat{2})
    &=f(\widehat{1})f(\widehat{2})\\
  &\Longrightarrow f(\widehat{3})=f(\widehat{1})\\
  &\overset{f\,\mathrm{inj.}}{\Longrightarrow}
      \widehat{3}=\widehat{1}\qquad\text{(fals).}
\end{aligned}
\]

\noindent La sfârșitul paginii mai este indicat monoidul de matrice
$(M_n(K),\cdot)$, cu $A,B,C\in M_n(K)$ astfel încât
\[
  AB=CB,\qquad A\neq C.
\]
; fără preambul sau pagină nouă.
% NOTE EDITORIALE:
% - Exercițiul 6 are două puncte. Demonstrația tratează separat
%   cazul unui monoid și cazul unui semigrup fără element neutru.
% - Sursa notează x*infinit = infinit*x = x, dar nu scrie separat
%   infinit*infinit = infinit; această egalitate este necesară pentru
%   ca operația să fie definită pe întregul M ca monoid.
% - Manuscrisul folosește generic Z_n. Egalitatea ulterioară
%   3-căciulă*2-căciulă = 1-căciulă*2-căciulă, cu 3 != 1,
%   este valabilă pentru n=4, nu pentru orice n. Păstrăm Z_n.
% - Exemplul final cu matrice este doar indicat în manuscris;
%   nu adăugăm o demonstrație care nu este scrisă acolo.
% - Prescurtările și textul foarte șters nu sunt reconstruite
%   dincolo de pasajele lizibile.

\subsection*{Seminar 4 - Pagina 2}

\noindent\underline{Exercițiul 6.}

\noindent\textbf{(a)} Să se arate că orice semigrup $S$ se poate
scufunda într-un monoid: există un monoid $M$ și un morfism injectiv
de semigrupuri
\[
  f:S\longrightarrow M,\qquad f(xy)=f(x)f(y).
\]

\noindent\textbf{(b)} Există și un monoid care nu se poate scufunda
într-un grup: există un monoid $M$ astfel încât nu există un grup $G$
și un morfism injectiv de monoizi $f:M\to G$.

\medskip
\noindent\underline{Demonstrație.}

\noindent\textbf{(a)} Fie $S$ un semigrup.

\smallskip
\noindent\textbf{I.} Dacă $S$ este monoid, luăm $M=S$ și
\[
  i:S\longrightarrow M=S,\qquad i(x)=x.
\]

\noindent\textbf{II.} Dacă $S$ este semigrup, dar nu este monoid,
fie $\infty\notin S$ și notăm
\[
  M\overset{\mathrm{not}}{=}S\cup\{\infty\}.
\]
Pentru $x,y\in M$, operația este definită în manuscris prin
\[
\begin{gathered}
  x\circ y\ \overset{\mathrm{def}}{=}\ xy
     \quad(x,y\in S),\\
  \infty\circ x\ \overset{\mathrm{def}}{=}\
     x\circ\infty\ \overset{\mathrm{def}}{=}\ x
     \quad(x\in S).
\end{gathered}
\]
Atunci $(M,\circ,\infty)$ este monoid. Fie
\[
  f:S\longrightarrow S\cup\{\infty\}=M,\qquad f(x)=x;
\]
$f$ este morfism injectiv de semigrupuri.

\medskip
\noindent\textbf{(b)} $(\mathbb Z_n,\cdot)$ este monoid, unde
\[
  \mathbb Z_n=\{\widehat{0},\widehat{1},\ldots,
       \widehat{n-1}\}.
\]
Presupunem că există un grup $G$ și un morfism injectiv de monoizi
$f:\mathbb Z_n\to G$, astfel încât
\[
  f(xy)=f(x)f(y),\qquad f(\widehat{1})=e.
\]
Argumentul prin simplificare din manuscris este
\[
\begin{aligned}
  f(\widehat{3})f(\widehat{2})
    &=f(\widehat{1})f(\widehat{2})\\
  &\Longrightarrow f(\widehat{3})=f(\widehat{1})\\
  &\overset{f\,\mathrm{inj.}}{\Longrightarrow}
      \widehat{3}=\widehat{1}\qquad\text{(fals).}
\end{aligned}
\]

\noindent La sfârșitul paginii mai este indicat monoidul de matrice
$(M_n(K),\cdot)$, cu $A,B,C\in M_n(K)$ astfel încât
\[
  AB=CB,\qquad A\neq C.
\]
